\chapter{KẾT LUẬN VÀ HƯỚNG PHÁT TRIỂN}

\section{Kết luận}

Qua đề tài, nhóm đã tự cài DBSCAN và Incremental DBSCAN bằng Python theo hướng đối tượng. DBSCAN dùng để gom cụm ban đầu; phần Incremental DBSCAN hỗ trợ thêm và xóa điểm, sử dụng lại thông tin lân cận và trạng thái cụm. Các bước cập nhật được xây dựng theo những ý chính trong tài liệu Incremental DBSCAN, gồm \(\mathrm{UpdSeed}_{Ins}\), \(\mathrm{UpdSeed}_{Del}\), kiểm tra kết nối và xử lý trường hợp cụm bị tách.

Nhóm đã minh họa các trường hợp INSERT và DELETE bằng ví dụ 2D, đồng thời chạy DBSCAN trên tọa độ ATM và khảo sát luồng thẻ theo các lát 2D, 3D, 4D. Silhouette do nhóm tự cài được dùng làm chỉ số chính để nhận xét chất lượng cụm; số cụm và tỷ lệ nhiễu được xem kèm để tránh kết luận chỉ từ một con số.

Trong file đối chiếu hiện có, 317 trạng thái sau khi fit hoặc cập nhật đều tương đương với DBSCAN chạy lại. Trên benchmark 20 cập nhật, DBSCAN scikit-learn fit lại nhanh hơn bản OOP ở cả 2D, 3D và 4D; HDBSCAN nhanh hơn ở 2D/3D và chậm hơn một chút ở 4D. OPTICS chậm nhất trong cả ba cấu hình. Bước fit D0 ban đầu của bản OOP cũng lâu hơn các bản thư viện. Vì vậy, thí nghiệm này kiểm tra được tính đúng và cho thấy thời gian phụ thuộc vào cách cài đặt, nhưng chưa chứng minh Incremental DBSCAN OOP có lợi thế tốc độ trên luồng đang đo.

\section{Hạn chế của đề tài}

Các dữ liệu đều là mô phỏng, nên chưa cho biết thuật toán hoạt động thế nào trên giao dịch thật. Benchmark thời gian chỉ chạy 20 cập nhật đầu trong luồng 300 cập nhật; các chiều 2D, 3D và 4D lấy từ cùng dữ liệu nên không phải các thử nghiệm độc lập.

Các hằng số chuẩn hóa của luồng thẻ được tính từ toàn bộ dữ liệu mô phỏng, bao gồm cả các cập nhật ở phía sau. Cách này giúp giữ tọa độ nhất quán giữa các thuật toán để so sánh, nhưng chưa mô phỏng một hệ thống online chỉ biết dữ liệu đến hiện tại.

Ba phương pháp scikit-learn được fit lại toàn bộ sau mỗi cập nhật, vì vậy đây là đối chiếu với cách phân cụm lại chứ không phải với một thư viện Incremental DBSCAN. Thời gian phụ thuộc vào môi trường chạy và cách cài đặt. Nhóm chưa đo bộ nhớ, và bộ tham số cố định trong benchmark cũng không được xem là tối ưu cho mọi dữ liệu.

\section{Hướng phát triển}

Nếu có thêm thời gian và dữ liệu phù hợp, nhóm có thể tiếp tục theo các hướng sau:

\begin{itemize}
    \item Thử trên dữ liệu đã ẩn thông tin cá nhân hoặc dữ liệu thực được phép sử dụng, rồi so sánh với kết quả trên dữ liệu mô phỏng.
    \item Phát lại nhiều bước cập nhật hơn và kiểm tra thêm các kích thước, số chiều để xem chi phí thay đổi ra sao.
    \item Tính tham số chuẩn hóa theo dữ liệu đã quan sát tại từng thời điểm để mô phỏng luồng online sát hơn.
    \item Đo thêm bộ nhớ và tìm cách giảm thời gian fit ban đầu, chẳng hạn bằng cấu trúc tra cứu lân cận phù hợp.
    \item Nếu tìm được một cài đặt Incremental DBSCAN có thể đối chiếu, so sánh thêm với cách cập nhật tăng dần đó thay vì chỉ so với các thuật toán fit lại toàn bộ.
\end{itemize}

Các hướng này là phần có thể làm tiếp, không nằm trong kết quả đã báo cáo ở đề tài hiện tại.
